% Options for packages loaded elsewhere
\PassOptionsToPackage{unicode}{hyperref}
\PassOptionsToPackage{hyphens}{url}
\PassOptionsToPackage{dvipsnames,svgnames,x11names}{xcolor}
\documentclass[
  11pt,
]{article}
\usepackage{xcolor}
\usepackage[margin=25mm]{geometry}
\usepackage{amsmath,amssymb}
\setcounter{secnumdepth}{-\maxdimen} % remove section numbering
\usepackage{iftex}
\ifPDFTeX
  \usepackage[T1]{fontenc}
  \usepackage[utf8]{inputenc}
  \usepackage{textcomp} % provide euro and other symbols
\else % if luatex or xetex
  \usepackage{unicode-math} % this also loads fontspec
  \defaultfontfeatures{Scale=MatchLowercase}
  \defaultfontfeatures[\rmfamily]{Ligatures=TeX,Scale=1}
\fi
\usepackage{lmodern}
\ifPDFTeX\else
  % xetex/luatex font selection
\fi
% Use upquote if available, for straight quotes in verbatim environments
\IfFileExists{upquote.sty}{\usepackage{upquote}}{}
\IfFileExists{microtype.sty}{% use microtype if available
  \usepackage[]{microtype}
  \UseMicrotypeSet[protrusion]{basicmath} % disable protrusion for tt fonts
}{}
\makeatletter
\@ifundefined{KOMAClassName}{% if non-KOMA class
  \IfFileExists{parskip.sty}{%
    \usepackage{parskip}
  }{% else
    \setlength{\parindent}{0pt}
    \setlength{\parskip}{6pt plus 2pt minus 1pt}}
}{% if KOMA class
  \KOMAoptions{parskip=half}}
\makeatother
\usepackage{listings}
\newcommand{\passthrough}[1]{#1}
\lstset{defaultdialect=[5.3]Lua}
\lstset{defaultdialect=[x86masm]Assembler}
\usepackage{longtable,booktabs,array}
\usepackage{calc} % for calculating minipage widths
% Correct order of tables after \paragraph or \subparagraph
\usepackage{etoolbox}
\makeatletter
\patchcmd\longtable{\par}{\if@noskipsec\mbox{}\fi\par}{}{}
\makeatother
% Allow footnotes in longtable head/foot
\IfFileExists{footnotehyper.sty}{\usepackage{footnotehyper}}{\usepackage{footnote}}
\makesavenoteenv{longtable}
\setlength{\emergencystretch}{3em} % prevent overfull lines
\providecommand{\tightlist}{%
  \setlength{\itemsep}{0pt}\setlength{\parskip}{0pt}}
\emergencystretch=3em
\AtBeginDocument{\Urlmuskip=0mu plus 2mu}
\lstset{breaklines=true,basicstyle=\small\ttfamily,columns=fullflexible,keepspaces=true}
\setlength{\tabcolsep}{4pt}
\usepackage{iftex}
\ifXeTeX
  \usepackage{newunicodechar}
  \makeatletter
  \newcommand{\DeclareUnicodeCharacter}[2]{%
    \begingroup
    \uccode`\~="#1\relax
    \uppercase{\def\temp{~}}%
    \expandafter\endgroup
    \expandafter\newunicodechar\expandafter{\temp}{#2}%
  }
  \makeatother
\fi
\ifdefined\DeclareUnicodeCharacter
\DeclareUnicodeCharacter{2192}{\ensuremath{\to}}
\DeclareUnicodeCharacter{2190}{\ensuremath{\leftarrow}}
\DeclareUnicodeCharacter{2194}{\ensuremath{\leftrightarrow}}
\DeclareUnicodeCharacter{21D2}{\ensuremath{\Rightarrow}}
\DeclareUnicodeCharacter{21A6}{\ensuremath{\mapsto}}
\DeclareUnicodeCharacter{2212}{\ensuremath{-}}
\DeclareUnicodeCharacter{2264}{\ensuremath{\leq}}
\DeclareUnicodeCharacter{2265}{\ensuremath{\geq}}
\DeclareUnicodeCharacter{2260}{\ensuremath{\neq}}
\DeclareUnicodeCharacter{2113}{\ensuremath{\ell}}
\DeclareUnicodeCharacter{03A3}{\ensuremath{\Sigma}}
\DeclareUnicodeCharacter{03A6}{\ensuremath{\Phi}}
\DeclareUnicodeCharacter{0393}{\ensuremath{\Gamma}}
\DeclareUnicodeCharacter{221E}{\ensuremath{\infty}}
\DeclareUnicodeCharacter{2208}{\ensuremath{\in}}
\DeclareUnicodeCharacter{2209}{\ensuremath{\notin}}
\DeclareUnicodeCharacter{211A}{\ensuremath{\mathbb{Q}}}
\DeclareUnicodeCharacter{211D}{\ensuremath{\mathbb{R}}}
\DeclareUnicodeCharacter{2124}{\ensuremath{\mathbb{Z}}}
\DeclareUnicodeCharacter{2115}{\ensuremath{\mathbb{N}}}
\DeclareUnicodeCharacter{2200}{\ensuremath{\forall}}
\DeclareUnicodeCharacter{2203}{\ensuremath{\exists}}
\DeclareUnicodeCharacter{2286}{\ensuremath{\subseteq}}
\DeclareUnicodeCharacter{00D7}{\ensuremath{\times}}
\DeclareUnicodeCharacter{03B4}{\ensuremath{\delta}}
\DeclareUnicodeCharacter{2202}{\ensuremath{\partial}}
\DeclareUnicodeCharacter{03B2}{\ensuremath{\beta}}
\DeclareUnicodeCharacter{03B1}{\ensuremath{\alpha}}
\DeclareUnicodeCharacter{03C6}{\ensuremath{\phi}}
\DeclareUnicodeCharacter{03B3}{\ensuremath{\gamma}}
\DeclareUnicodeCharacter{00B9}{\ensuremath{^1}}
\DeclareUnicodeCharacter{00B2}{\ensuremath{^2}}
\DeclareUnicodeCharacter{2070}{\ensuremath{^0}}
\DeclareUnicodeCharacter{2080}{\ensuremath{_0}}
\DeclareUnicodeCharacter{2081}{\ensuremath{_1}}
\DeclareUnicodeCharacter{2082}{\ensuremath{_2}}
\fi
\usepackage{bookmark}
\IfFileExists{xurl.sty}{\usepackage{xurl}}{} % add URL line breaks if available
\urlstyle{same}
\hypersetup{
  pdftitle={Reviewing claims against code and evidence},
  pdfauthor={Jeremy H. Carroll},
  colorlinks=true,
  linkcolor={Maroon},
  filecolor={Maroon},
  citecolor={Blue},
  urlcolor={Blue},
  pdfcreator={LaTeX via pandoc}}

\title{Reviewing claims against code and evidence}
\author{Jeremy H. Carroll}
\date{2026-10-02 · Version 0.1.0 · Review protocol · Author-review
revision}

\begin{document}
\makeatletter
\begin{center}
{\LARGE \@title\par}
\vspace{0.8em}
{\normalsize \@author\par}
\vspace{0.5em}
{\small \@date\par}
\end{center}
\makeatother

\textbf{Code baseline:}
\nolinkurl{52cb5457f6c485ba548cfd292f6012a941ed22cf}

\subsection{Purpose}\label{purpose}

This protocol helps a reviewer decide what a technical claim is
supported by. It is a companion to the Proposal Fidelity Protocol (PFP)
papers, operator guide, and curriculum guide. ``Adversarial'' means
deliberately testing plausible ways the claim could fail. It does not
require hostile language, confidence theater, or agreement among
multiple models.

The unit of review is a claim together with its inputs, implementation,
and evidence. A successful check supports the property it checked under
its assumptions. A failed check or missing prerequisite is useful
information to preserve. This document provides a review procedure and
worked examples. The only recorded model run it uses is a retained
ten-trial trace, which Section 4.1 recounts as an example of reading a
reported catch rate. Appendix A applies the procedure to one gate in
this repository, the terminal gate of the seven-act Adversary engine, with
four falsifiers replayed on the audited code. It reports no measured
psychological effect, search-escape rate, or general automated truth
detector.

\subsection{1. Make the claim small enough to
inspect}\label{make-the-claim-small-enough-to-inspect}

Replace a broad sentence such as ``the system verifies repairs'' with a
concrete statement:

\begin{quote}
For this declared triangle, a replacement proposal that increases the
checked quotient minimum is refused, and the active claims remain
unchanged.
\end{quote}

The reviewer can now identify the graph, original claims, replacement,
score definition, relevant branch, and expected stored state. If one is
missing, ask the author to supply it before evaluating the broad claim.

An \textbf{assumption} is a condition required by the statement. A
\textbf{witness} is supplied data used to establish a particular
property, such as a candidate potential or dual vector. A
\textbf{receipt} is an execution record. An \textbf{open obligation} is
a specific unanswered question needed for a stronger conclusion.

Keep a stable source revision and exact artifact paths. When a source
file changes, attach a new observation to the new bytes rather than
silently reusing the old result.

\subsection{2. Inspect the connection between language and
behavior}\label{inspect-the-connection-between-language-and-behavior}

For each important term, make a three-column note:

\begin{longtable}[]{@{}
  >{\raggedright\arraybackslash}p{(\linewidth - 4\tabcolsep) * \real{0.3333}}
  >{\raggedright\arraybackslash}p{(\linewidth - 4\tabcolsep) * \real{0.3333}}
  >{\raggedright\arraybackslash}p{(\linewidth - 4\tabcolsep) * \real{0.3333}}@{}}
\toprule\noalign{}
\begin{minipage}[b]{\linewidth}\raggedright
Name in the claim
\end{minipage} & \begin{minipage}[b]{\linewidth}\raggedright
Property intended
\end{minipage} & \begin{minipage}[b]{\linewidth}\raggedright
Operation actually performed
\end{minipage} \\
\midrule\noalign{}
\endhead
\bottomrule\noalign{}
\endlastfoot
``exact repair'' & Changed claims satisfy the declared graph relation &
A native result is checked under the selected graph inputs \\
``notebook grading'' & A learner's answer is assessed & Perhaps only a
supplied \passthrough{\lstinline!execution\_status!} string is
counted \\
``independent verifier'' & Errors can be detected without trusting the
proposal & Perhaps only two parameter-hash fields are compared \\
``replay'' & A recorded result is recomputed from its dependencies &
Perhaps only one stored digest is compared with another \\
\end{longtable}

The third column comes from inspecting the actual call path. Matching
names or comments are insufficient. Conversely, a deliberately small
check is useful when its stated property is equally small.

The package's
\href{https://github.com/dbccpage/the_symbol_game/blob/52cb5457f6c485ba548cfd292f6012a941ed22cf/atlas/coda_atlas/symbol_game/009_pfp_adversarial_dialectics_and_terminal_reality_gates/code/tests/test_009_adversarial_reality_gate.py}{semiotic-audit
fixture} compares the strings \passthrough{\lstinline!signified!} and
\passthrough{\lstinline!interpretant!}. It illustrates mismatched
descriptions. It does not inspect downstream code, determine
mathematical equivalence, or audit every execution path. Use a reviewer
to establish those connections unless a separately validated tool
actually performs the work.

\subsection{3. Work through a concrete false
reading}\label{work-through-a-concrete-false-reading}

Take the triangle claims \(h=(1/3,1/3,1/3)\). Their cycle sum is one.
Suppose an author says that subtracting endpoint-potential differences
``removes the cycle obstruction.''

Choose potentials \(f=(0,1/3,2/3)\). Their edge differences are
\(df=(1/3,1/3,-2/3)\), so \(h-df=(0,0,1)\). The discrepancy has been
concentrated, but the cycle sum remains one. Every potential produces
edge differences summing to zero, so the objection concerns the proposed
mechanism, not an unlucky choice of numbers.

A different operation can replace the claims by \(df\), whose cycle sum
is zero. That statement is supportable because the declared input
changed. The reviewer should request this distinction in the main
explanation, not accept a caution several pages later that contradicts
the headline.

The
\href{https://github.com/dbccpage/the_symbol_game/blob/52cb5457f6c485ba548cfd292f6012a941ed22cf/education/the_symbol_game/lessons/system_integration/100_novices/04_cycle_repair.ipynb}{cycle
lesson} reproduces the first reading with these values, and the
\href{../005/paper.pdf}{operator guide} shows both readings on an
integer triangle. This package's
\href{https://github.com/dbccpage/the_symbol_game/blob/52cb5457f6c485ba548cfd292f6012a941ed22cf/atlas/coda_atlas/symbol_game/009_pfp_adversarial_dialectics_and_terminal_reality_gates/jupyter_notebooks/009_triangle_cycle_sum.ipynb}{companion
notebook} and
\href{https://github.com/dbccpage/the_symbol_game/blob/52cb5457f6c485ba548cfd292f6012a941ed22cf/atlas/coda_atlas/symbol_game/009_pfp_adversarial_dialectics_and_terminal_reality_gates/code/tests/test_triangle_cycle_sum.py}{triangle
tests} recompute the numbers above with exact fractions.

\subsection{4. Design a probe that could change the
verdict}\label{design-a-probe-that-could-change-the-verdict}

A \textbf{positive control} is a case expected to pass. A
\textbf{negative control} is a case with a known defect expected to
fail. The pair tests whether a checker distinguishes cases, rather than
accepting or refusing everything.

A \textbf{linear program (LP)} minimizes a linear objective subject to
linear constraints. A feasible \textbf{primal witness} obeys the
original constraints; a feasible \textbf{dual witness} obeys the
associated lower-bound constraints. Matching objective values certify
optimality. In the graph example, dual closure means balanced flow at
every vertex; the \href{../001/paper.pdf}{finite graph paper} gives the
argument.

For an exact LP certificate, first supply feasible primal and dual
witnesses with equal objectives. Then independently corrupt a primal
bound, break dual closure, and retain a feasible but nonmatching dual
objective. The
\href{https://github.com/dbccpage/the_symbol_game/blob/52cb5457f6c485ba548cfd292f6012a941ed22cf/education/the_symbol_game/lessons/system_integration/200_intermediates/02_exact_lp_certificates.ipynb}{LP
lesson} supplies these three distinct failure cases. These are
source-inspected cases; no fresh replay result is reported here. Equal
objective values alone cannot replace feasibility.

For a digest-bound packet, change node content while keeping its old
digest key. The
\href{https://github.com/dbccpage/the_symbol_game/blob/52cb5457f6c485ba548cfd292f6012a941ed22cf/education/the_symbol_game/lessons/system_integration/100_novices/01_codecs_and_tamper_evidence.ipynb}{codec
lesson} contains that mutation. Checking that a string has 64 characters
would not test the same property.

Record the original and changed input, expected distinction, observed
result, and reason the probe targets the claim. A probe passing is
evidence for that case; exhaustiveness requires an independently
justified finite enumeration or proof.

\subsubsection{4.1 Reading a reported catch
rate}\label{reading-a-reported-catch-rate}

The same controls apply to a reported detection rate. Paper 001's
package retains one paired run in which a proposer model classified
small Hubbard-model instances as sign-free or as carrying an open sign
obstruction, and an adversary model challenged or accepted each
classification. The run covers two frustrated lattices, a five-site ring
and a six-site cluster of corner-sharing triangles, each at five
interaction strengths, for 10 trials. Its
\href{https://github.com/dbccpage/the_symbol_game/blob/52cb5457f6c485ba548cfd292f6012a941ed22cf/atlas/coda_atlas/symbol_game/001_pfp_discrete_cochain_retraction_and_toy_involutions/results/demo2_r7_adversarial_debate_receipt.json}{replay
receipt} reports the proposer as
\nolinkurl{meta-llama/llama-3.1-8b-instruct} and the adversary as
\passthrough{\lstinline!deepseek/deepseek-chat!}. The
\href{https://github.com/dbccpage/the_symbol_game/blob/52cb5457f6c485ba548cfd292f6012a941ed22cf/atlas/coda_atlas/symbol_game/001_pfp_discrete_cochain_retraction_and_toy_involutions/results/demo2_r7_adversarial_debate_1790210039.jsonl}{trace}
itself records no model identifiers.

Each trial's reference label comes from a program check, not from a
model. Six proposals claimed a sign-free certificate and were refused by
Paper 001's certificate checker; they are labeled unsound. The other
four claimed an open sign obstruction and were not refused at the time;
they are labeled sound. The receipt describes these as historical
verifier labels, not independent physical ground truth.

\begin{longtable}[]{@{}llll@{}}
\toprule\noalign{}
Checker label & Adversary challenged & Adversary concurred & Trials \\
\midrule\noalign{}
\endhead
\bottomrule\noalign{}
\endlastfoot
Unsound (refused) & 6 & 0 & 6 \\
Sound (not refused) & 2 & 2 & 4 \\
All trials & 8 & 2 & 10 \\
\end{longtable}

In these 10 trials the adversary challenged all six unsound proposals
and missed none. It also challenged two of the four sound-labeled
proposals and passed the other two. In both of those challenges its
counter-proposal was the same open-obstruction claim, so it disputed the
reasoning rather than the conclusion. Counting the trace field
\passthrough{\lstinline!adversary\_caught\_am!} instead gives 10 of 10,
because that field is true in every record, including both concurrences.
Its name suggests detection; its values do not measure it.

Four observations limit what the run shows:

\begin{itemize}
\tightlist
\item
  All ten trials ended in the same recorded terminal state, an open sign
  obstruction, whether the adversary challenged or concurred.
\item
  The six refusals come from the checker. The adversary's challenges
  agreed with them; the run does not show the adversary stopping
  anything the checker would have accepted.
\item
  A later replay of the ten recorded proposals through a revised version
  of the checker, stored in the same receipt, refuses all ten. The four
  obstruction claims fail because none supplies a checked obstruction
  witness. The sound label therefore means only that the checker did not
  refuse the claim at the time.
\item
  The run says nothing about other models, lattices, or prompts.
\end{itemize}

The package test
\href{https://github.com/dbccpage/the_symbol_game/blob/52cb5457f6c485ba548cfd292f6012a941ed22cf/atlas/coda_atlas/symbol_game/009_pfp_adversarial_dialectics_and_terminal_reality_gates/code/tests/test_recorded_bout_counts.py}{\nolinkurl{test_recorded_bout_counts.py}}
recounts these cells from the trace and checks the receipt. It needs the
full repository, because both files are kept in Paper 001's package, and
it skips when this package is used on its own.

\subsection{5. Read formal statements at their actual
scope}\label{read-formal-statements-at-their-actual-scope}

Start with the theorem's definitions, hypotheses, and conclusion. Then
identify the bridge from the paper's objects to those formal objects.

This package's \passthrough{\lstinline!lean4!} folder holds copies of
two of the files discussed below, together with the two modules that the
second file imports. Each copy matches the repository's
\passthrough{\lstinline!lean4!} file at the code baseline except for a
two-line provenance comment at the top. They build with Lean 4.34.0 and
no external libraries. The
\href{https://github.com/dbccpage/the_symbol_game/blob/52cb5457f6c485ba548cfd292f6012a941ed22cf/atlas/coda_atlas/symbol_game/009_pfp_adversarial_dialectics_and_terminal_reality_gates/lean4/axiom_report.txt}{axiom
report} lists every theorem in the four copied modules; none uses an
axiom other than \passthrough{\lstinline!propext!} and
\passthrough{\lstinline!Quot.sound!}.

\href{https://github.com/dbccpage/the_symbol_game/blob/52cb5457f6c485ba548cfd292f6012a941ed22cf/atlas/coda_atlas/symbol_game/009_pfp_adversarial_dialectics_and_terminal_reality_gates/lean4/Level15AuthorityCalculus.lean}{\mbox{\nolinkurl{Level15AuthorityCalculus.lean}}}
defines distinct state constructors and formalizes that designated
proposal constructors differ from the terminal-authority constructor. It
also defines an ``independent verifier'' by unequal parameter hashes and
shows that a state does not meet that definition against itself. These
facts do not establish statistical independence of different models or
guarantee that a runtime preserves the state partition.

\href{https://github.com/dbccpage/the_symbol_game/blob/52cb5457f6c485ba548cfd292f6012a941ed22cf/atlas/coda_atlas/symbol_game/009_pfp_adversarial_dialectics_and_terminal_reality_gates/lean4/TriggerNotLicense.lean}{\mbox{\nolinkurl{TriggerNotLicense.lean}}}
constructs a finite case with an active trigger and zero lift budget. It
shows that a trigger can exist without a licensed target under those
definitions. It does not preclude independently licensed actions under
alternate assumptions.

The repository's terminal-routing source
\href{https://github.com/dbccpage/the_symbol_game/blob/52cb5457f6c485ba548cfd292f6012a941ed22cf/lean4/TerminalizerPolicy.lean}{\passthrough{\lstinline!TerminalizerPolicy!}}
imports a Mathlib module and carries a source warning that its
declarations were not typechecked against the named package
configuration. It is not copied into this package. Its table should
therefore not be advertised here as a freshly compiled termination
theorem. A build receipt must name the actual toolchain, dependencies,
targets, and result.

These examples make the review rule practical: cite what the formal
statement says, and identify any missing implementation or modeling
correspondence.

\subsection{6. State the review outcome without promoting
it}\label{state-the-review-outcome-without-promoting-it}

Use the following report terms as a declared review convention. They are
not a new runtime API.

\begin{longtable}[]{@{}
  >{\raggedright\arraybackslash}p{(\linewidth - 2\tabcolsep) * \real{0.5000}}
  >{\raggedright\arraybackslash}p{(\linewidth - 2\tabcolsep) * \real{0.5000}}@{}}
\toprule\noalign{}
\begin{minipage}[b]{\linewidth}\raggedright
Review disposition
\end{minipage} & \begin{minipage}[b]{\linewidth}\raggedright
Meaning
\end{minipage} \\
\midrule\noalign{}
\endhead
\bottomrule\noalign{}
\endlastfoot
Supported within scope & The inspected evidence supports the stated
claim under recorded assumptions \\
Refused claim & A counterexample, failed required check, or
contradiction defeats the stated claim \\
Open obligation & A particular missing bridge or result prevents the
stronger conclusion \\
Not run & The proposed execution has not been performed in this
review \\
\end{longtable}

A package may contain all four dispositions on different claims. Report
them individually. Forwarding a supported result for another review is a
workflow step; it grants no publication, deployment, or domain
authority.

The package's
\href{https://github.com/dbccpage/the_symbol_game/blob/52cb5457f6c485ba548cfd292f6012a941ed22cf/atlas/coda_atlas/symbol_game/009_pfp_adversarial_dialectics_and_terminal_reality_gates/code/tests/test_009_adversarial_reality_gate.py}{toy
verdict function} maps supplied rational scores and flags to
\passthrough{\lstinline!PASS\_TO\_MASTER!},
\passthrough{\lstinline!REFUSE!}, or
\passthrough{\lstinline!OPEN\_RESIDUAL!}. It does not compute those
scores or establish that the flags describe reality. Its supplied-input
tests cannot support a claim of comprehensive semantic verification.

\subsection{7. Preserve enough information for another
reviewer}\label{preserve-enough-information-for-another-reviewer}

A compact record should contain the claim text, source revision, file
path, assumptions, exact input, command, environment or executable pin,
observed outcome, and remaining obligation. Include the failing case and
any skipped dependency. Identify whether a stored output was only read
or was recomputed.

For example, a review request can be recorded before execution:

\begin{lstlisting}
claim_id: triangle-replacement-01
source_revision: 52cb5457f6c485ba548cfd292f6012a941ed22cf
original_claims: ["1", "1", "1"]
proposed_claims: ["5", "0", "0"]
expected_status: DefectIncrease
run_status: NOT_RUN
next_check: Run the focused native-backed session test and inspect stored state
decision_scope: review-only
\end{lstlisting}

This is a template, not an execution receipt. Replace the pending fields
only with observed evidence. A digest binds a declared byte sequence;
authenticating a producer and evaluating content are separate tasks.
When claiming canonical JSON, test the actual requirements of
\href{https://www.rfc-editor.org/rfc/rfc8785}{RFC 8785}, including
primitive serialization, rather than equating key sorting with complete
conformance.

\subsection{8. Handle stalled work
explicitly}\label{handle-stalled-work-explicitly}

When several attempts produce no useful change, record the unchanged
residual and the attempted move family. Check whether the current
representation permits the requested result. The triangle above
demonstrates why more potential adjustments cannot change its cycle sum.

A new representation or new allowed move is a proposal that needs its
own assumptions and checks. A stagnation trigger can suggest review; it
does not license changing protected state or guarantee an escape. The
local fixture only checks a supplied stagnation count, a locked-gauge
flag, and a positive supplied score. It does not implement an effective
search strategy.

\subsection{Related work and scope}\label{related-work-and-scope}

Necula's \href{https://doi.org/10.1145/263699.263712}{proof-carrying
code} provides a primary example of a producer supplying evidence that a
host checks against a declared policy. Saltzer and Schroeder's
\href{https://doi.org/10.1109/PROC.1975.9939}{protection principles}
provide security-design context for explicit permission and mediation.
This protocol borrows a review discipline from these works; it does not
inherit their technical results.

The useful deliverable is a review whose reader can trace each
conclusion to a precise statement, source, probe, or formal obligation.
More ambitious claims about automated reasoning effectiveness require
measurements or proofs beyond that trace.

\subsection{Limitations}\label{limitations}

This is a review procedure, not a measured method; no study of reviewer
accuracy or effort stands behind it. Section 4.1 recounts one retained
run of ten trials whose labels come from a program check, and it
supports no rate for other settings. Appendix A audits one gate at one
commit with four constructed inputs; it shows that each defect exists,
not how often it changes a verdict in use. The package fixtures decide
only from supplied values. The copied Lean modules prove narrow
statements about their own definitions, and the terminal-routing source
discussed in Section 5 states that it was not typechecked. This paper
assigns no authorization and changes no jurisdiction. It carries no
operational authority or institutional standing.

\subsection{Verification for this
revision}\label{verification-for-this-revision}

The package verifier (\passthrough{\lstinline!verify\_deposit.py!})
checks the manifest and the two-way file census, and
\passthrough{\lstinline!tools/rebuild\_release.py --check!} checks the
release archive against the manifests. The package tests, lint, Lean
build, and notebook replay are executed by the dated review receipt
recorded with this revision; that receipt lists each command, tool
version, exit status, and tested scope. Historical run records cited
here, such as the Section 4.1 trace in Paper 001's package, keep their
original dates and scope and are not current execution evidence.

\subsection{References}\label{references}

\begin{itemize}
\tightlist
\item
  Necula, G. C. (1997).
  \href{https://doi.org/10.1145/263699.263712}{Proof-Carrying Code}. In
  \emph{Proceedings of the 24th ACM SIGPLAN-SIGACT Symposium on
  Principles of Programming Languages (POPL '97)}, pp.~106--119. ACM.
  \nolinkurl{doi:10.1145/263699.263712.}
\item
  Saltzer, J. H., and Schroeder, M. D. (1975).
  \href{https://doi.org/10.1109/PROC.1975.9939}{The Protection of
  Information in Computer Systems}. \emph{Proceedings of the IEEE},
  63(9), 1278--1308. \nolinkurl{doi:10.1109/PROC.1975.9939.}
  \href{https://web.mit.edu/Saltzer/www/publications/protection/}{Author-hosted
  copy}.
\item
  Rundgren, A., Jordan, B., and Erdtman, S. (2020).
  \href{https://www.rfc-editor.org/rfc/rfc8785}{JSON Canonicalization
  Scheme (JCS)}. RFC 8785, RFC Editor, June 2020. Independent Submission
  stream, Informational. doi:10.17487/RFC8785.
\end{itemize}

\subsection{Appendix A. Where a gate's verdict comes
from}\label{appendix-a.-where-a-gates-verdict-comes-from}

This appendix applies Sections 2, 4 and 7 to one gate in this
repository: the terminal gate of the seven-act Adversary engine
(\passthrough{\lstinline!omega\_engine\_adversary!}) at commit
\nolinkurl{30609de2c3928238715375a9bac608256fa381d8}. The gate is meant
to authorize a proposal only when exact defect measures of the input are
zero: a quotient defect, written \(\Phi_1\), and a cycle defect, written
\(\Gamma_2\). Otherwise it should refuse, or record a bounded residual.
The review question is the one of Section 2: for each verdict, which
operation actually produced it?

Four falsifiers answer that question. Each fixes an input, states the
status a sound gate must give, and records the status the audited code
gave. They were run on the audited commit by a script of this
collection, \nolinkurl{tools/replay_adversary_gate_falsifiers.py}, which
extracts the needed packages with \passthrough{\lstinline!git archive!}
and calls single act functions with a stub runtime that has no event
sinks and no model port. Nothing is written and no model is called. The
observed values below are those of its record,
\nolinkurl{verification/adversary_gate_falsifiers.json} (both paths
relative to the collection folder). File and line references are to the
audited commit.

Table A.1. The name of each operation and what it does.

\begin{longtable}[]{@{}
  >{\raggedright\arraybackslash}p{(\linewidth - 4\tabcolsep) * \real{0.3333}}
  >{\raggedright\arraybackslash}p{(\linewidth - 4\tabcolsep) * \real{0.3333}}
  >{\raggedright\arraybackslash}p{(\linewidth - 4\tabcolsep) * \real{0.3333}}@{}}
\toprule\noalign{}
\begin{minipage}[b]{\linewidth}\raggedright
Name in the gate
\end{minipage} & \begin{minipage}[b]{\linewidth}\raggedright
Property intended
\end{minipage} & \begin{minipage}[b]{\linewidth}\raggedright
Operation actually performed
\end{minipage} \\
\midrule\noalign{}
\endhead
\bottomrule\noalign{}
\endlastfoot
``Gauge break'' and ``carrier lift'' (Act IV) & A lifted carrier on
which the input's obstruction can be examined & A node is added and the
cochain is replaced by zero on every edge
(\nolinkurl{adversary/use_cases/act_helpers.py}, lines 480 to
498); Act II then scores the zeros
(\nolinkurl{use_cases/bout_service.py}, lines 266 to 277) \\
``Independent witness'' and ``countermodel'' (Act V) & A critique blocks
a candidate when it establishes a defect & A critique blocks when its
text contains the word ``DEFECT'' or ``FLAW''
(\passthrough{\lstinline!act\_helpers.py!}, line 642) \\
``Evidence references'' and ``obstruction non-vanishing receipt'' (Act
VII) & Each reference names stored bytes that support the verdict & With
no critique, the reference is the SHA-256 of the run key; the receipt's
identifier is the SHA-256 of a label string, and its fields
\passthrough{\lstinline!annihilates\_coboundary!} and
\passthrough{\lstinline!norm\_bounded!} are the constant
\passthrough{\lstinline!True!} (lines 917 to 932 and 1008 to 1015);
every event record's \passthrough{\lstinline!carrier\_hash!} is the
SHA-256 of the fixed bytes \passthrough{\lstinline!carrier!} (six sites,
from line 179) \\
``Novella'' and ``DDC tape'' of a challenge gate & A narrative and a
step record of what the run did & Fixed text and fixed steps into which
only the verdict, the counts and two numbers are substituted (C012 copy,
\nolinkurl{governance/adversary_gate.py}, lines 69 to 195) \\
\end{longtable}

Table A.2. Falsifiers: fixed input, expected status, observed status.

\begin{longtable}[]{@{}
  >{\raggedright\arraybackslash}p{(\linewidth - 6\tabcolsep) * \real{0.2500}}
  >{\raggedright\arraybackslash}p{(\linewidth - 6\tabcolsep) * \real{0.2500}}
  >{\raggedright\arraybackslash}p{(\linewidth - 6\tabcolsep) * \real{0.2500}}
  >{\raggedright\arraybackslash}p{(\linewidth - 6\tabcolsep) * \real{0.2500}}@{}}
\toprule\noalign{}
\begin{minipage}[b]{\linewidth}\raggedright
ID
\end{minipage} & \begin{minipage}[b]{\linewidth}\raggedright
Fixed input
\end{minipage} & \begin{minipage}[b]{\linewidth}\raggedright
A sound gate gives
\end{minipage} & \begin{minipage}[b]{\linewidth}\raggedright
The audited gate gave
\end{minipage} \\
\midrule\noalign{}
\endhead
\bottomrule\noalign{}
\endlastfoot
F1 & Triangle with edges (0,1), (1,2), (2,0) and cochain (1, 1, 1),
whose cycle sum is 3; default scorer; one candidate; no critique & Not
\passthrough{\lstinline!AUTHORIZED!}: no potential changes a cycle sum
(Section 3) & \passthrough{\lstinline!AUTHORIZED!} \\
F2a & Critic reply: ``Counterexample: the residual (0,0,3) has cycle sum
3, so the claimed repair is false.'' & Blocking & Not blocking; the
candidate survives \\
F2b & Critic reply: ``No defect found; the certificate checks.'' & Not
blocking & Blocking; the candidate is removed \\
F3 & Act VII with no critique, \(\Phi_1 = 1\), \(\Gamma_2 = 0\), and a
dual witness of 1 on edge (0,1) only & Each reference is the digest of
stored evidence; stated witness properties are checked & First reference
is the SHA-256 of the run key; the receipt states that the witness
annihilates coboundaries and is norm-bounded, although it is not
balanced \\
F4 & C012 gate with verdict \passthrough{\lstinline!REFUSE!}, 0 cases, 0
passed, 5 defects & No statement of repair, agreement or budget that the
inputs do not support & The novella states ``100\% concordance across
all cases'', a respected budget, and ``zero duality gap on all 46 judge
cases'', and names challenge C011; the step record contains an exact
coboundary repair with status ``repaired'' \\
\end{longtable}

\subsubsection{A.1 Input replacement
(F1)}\label{a.1-input-replacement-f1}

On the triangle, the default scorer reports 3 in place of \(\Phi_1\),
and \(\Gamma_2 = 0\). That scorer returns the \(\ell^1\) mass of the
cochain, not the quotient defect, and it returns \(\Gamma_2 = 0\) for
every input. Act IV fires because the score is positive and only one
candidate exists. It returns a lifted carrier whose cochain is zero on
every edge, and Act II on that cochain reports zero for both measures.
Act VII then authorizes. The verdict comes from the replaced input, not
from the proposal. This is the false reading of Section 3 in executable
form: replacing the claims changes the input, and the gate then
certifies the replacement.

Fix: commit \passthrough{\lstinline!a054c8ce1!} on branch
\nolinkurl{claude/mssn-study-corrections}. An Act IV lift now carries
the input cochain to the lifted carrier, and a lift that lowers the
measured obstruction is refused with
\nolinkurl{ADVERSARY_LIFT_LOWERED_OBSTRUCTION}. The default scorer reports
that nothing was measured, which cannot authorize, and Act VII requires
exact rational zeros. The laundering falsifiers are in
\nolinkurl{adversary/tests/test_honest_gate_falsifiers.py}.

\subsubsection{A.2 Keyword adjudication (F2a,
F2b)}\label{a.2-keyword-adjudication-f2a-f2b}

The critic's reply is a proposal like any other, and Section 2 asks what
the gate does with it. It searches the text for two words. A reply that
states a counterexample without either word does not block, and a reply
that says no defect was found does block. The status of the candidate
therefore depends on the critic's wording, not on a check of what the
critique claims.

Fix: commit \passthrough{\lstinline!a054c8ce1!}. A critique is now a
typed JSON block, and a critique that cannot be parsed blocks; no
keyword decides. The keyword falsifiers are in the same test file and in
\nolinkurl{adversary/tests/test_adversarial_m1_it2_1_challenger.py}.

\subsubsection{A.3 Phantom evidence references
(F3)}\label{a.3-phantom-evidence-references-f3}

An evidence reference is useful only if another reviewer can retrieve
the bytes it names (Section 7). With no critique, the audited gate
substitutes the SHA-256 of the run key, which names no stored evidence.
The identifier of the obstruction receipt is the SHA-256 of a label
built from the run key and the dual value, not the digest of the
receipt. The receipt itself asserts that the dual witness annihilates
coboundaries and is norm-bounded, but those fields are constants: the
witness of F3 has nonzero net flow at vertices 0 and 1, and no weights
were supplied to bound it. The event records carry the same
\passthrough{\lstinline!carrier\_hash!} in every run.

Fix: commit \passthrough{\lstinline!a054c8ce1!}. The obstruction and
deficit receipts are stored with the run under
\passthrough{\lstinline!evidence\_receipts!}, keyed by the digest of
their canonical bytes, and the verdict cites them by that digest. The
constant witness fields are removed, and the run's
\passthrough{\lstinline!receipt\_id!} is computed from its content.

\subsubsection{A.4 Templated receipts
(F4)}\label{a.4-templated-receipts-f4}

The challenge packages carry their own copies of a gate module that
writes a narrative (``novella'') and a hash-chained step record for each
run. In the C012 copy, both are templates. For a refused run with no
cases, the narrative still reports full agreement of the two checkers, a
respected budget and models under 70B, and its default summary claims
verified repairs on 46 judge cases; the step record lists an exact
simplex descent and a repaired coboundary. The copy also names challenge
C011, from which it was copied. The chain of digests over these steps is
internally consistent (Section 7), so a reader who checks only the chain
would accept them. The challenges folder holds 36 files named
\nolinkurl{governance/adversary_gate.py}, and the sentence about full
agreement occurs verbatim in four of them.

Fix: commit \passthrough{\lstinline!e6fbc15cf!} on the same branch. The
36 challenge copies of the gate now report a reference run as
\nolinkurl{REFERENCE_SELF_CONSISTENT} or
\nolinkurl{REFERENCE_SELF_INCONSISTENT} with a stated
\passthrough{\lstinline!gate\_scope!}, and the narrative is a table of
the measured counts with no templated claims. The tests are in
\nolinkurl{projects/mssn_sn_symgam/challenges/harness_tests/test_harness_honesty.py}.

\subsubsection{A.5 What the audit
supports}\label{a.5-what-the-audit-supports}

The four falsifiers show that, at the audited commit, an
\passthrough{\lstinline!AUTHORIZED!} status from this gate does not
establish that the proposal repaired its input, that a blocking critique
established a defect, or that a cited reference names evidence. They do
not show how often these defects changed a verdict in recorded runs,
which this appendix does not examine. In the corrected gate an
\passthrough{\lstinline!AUTHORIZED!} status requires an
\passthrough{\lstinline!ACCEPTED!} verdict from a candidate checker
bound to the candidate by digest: that verdict alone when no carrier is
declared, and that verdict together with exact zeros when one is. The
receipt field \passthrough{\lstinline!authorization\_evidence!} records
which of the two applied. The verdict therefore comes from a bound
checker receipt, the property Section 2 asks a reviewer to locate. On
the corrected commit the two Adversary test files named above pass (66 and
44 tests; the commands are in the collection's verification report).

\end{document}
